\begin{table}[!ht]
\centering
\caption{O que a categoria agregada esconde: penalidade racial por cor, setor, idade e
corte de renda. Penalidade salarial: HLM~M3 (mesmo bairro, escolaridade, idade, sexo,
jornada e estado), em \% de renda; por idade, M3 com a interação negro $\times$ faixa
etária. Razões de chance: GLMM de acesso e de teto de vidro (top 10\%), contra brancos
de mesmo perfil e bairro. Preço: parcela não explicada da Oaxaca--Blinder (A), contra
brancos. Setor público: empregados do setor público, inclusive militares e estatutários
(VD4009); privado: os demais ocupados. Cada recorte foi estimado em todos os seus
indivíduos, sem amostragem, com a especificação do modelo de referência.}
\label{tab:heterogeneidade}
\begin{tabular}{lrrrr}
\toprule
Recorte & Penalidade salarial (\%) & OR cargo qualificado & OR top 10\% & Preço (\%) \\
\midrule
Negro (categoria do trabalho) & 6,1 & 0,814 & 0,761 & 17,9 \\
\quad Pardo & 5,8 & 0,820 & 0,768 & 17,5 \\
\quad Preto & 7,4 & 0,785 & 0,725 & 20,0 \\
Setor privado & 7,2 & 0,783 & 0,685 & --- \\
Setor público & 4,6 & 0,786 & 0,843 & --- \\
Idade 14--29 anos & 1,7 & --- & --- & --- \\
Idade 30--39 anos & 5,8 & --- & --- & --- \\
Idade 40--49 anos & 6,9 & --- & --- & --- \\
Idade 50--64 anos & 8,8 & --- & --- & --- \\
Idade 65 anos ou mais & 14,7 & --- & --- & --- \\
Top 10\% dentro da UF & --- & --- & 0,762 & --- \\
\bottomrule
\end{tabular}
\par\smallskip
\footnotesize\emph{Como ler:} OR abaixo de~1 é menor chance que a de um branco
comparável; ``Preço'' é a parte do gap que as características observáveis não explicam.
Na última linha, o topo é definido dentro de cada estado, e não pelo corte nacional.
\end{table}
